\subsection{例 3：情形 L = A}

若 L = A，令 \(u(1) = x \in \mathbf{A}\) 。复形 \(\mathbf{K}(u)\) 的长度为 1，且我们有 \(\mathbf{K}_{0}(u) = \mathbf{K}_{1}(u) = \mathbf{A}\) 和 \(d_{1}(a) = xa\) ，因此对每个 A-模 M， \(\mathbf{K}_{0}(u, \mathbf{M})\) , \(\mathbf{K}_{1}(u, \mathbf{M})\) , \(\mathbf{K}^{0}(u, \mathbf{M})\) 和 \(\mathbf{K}^{1}(u, \mathbf{M})\) 与 M 等同，微分

\[
d _ {1}: \mathbf {K} _ {1} (u, \mathrm{M}) \rightarrow \mathbf {K} _ {0} (u, \mathrm{M}) \quad \text { and }\quad d ^ {0}: \mathbf {K} ^ {0} (u, \mathrm{M}) \rightarrow \mathbf {K} ^ {1} (u, \mathrm{M})
\]

是 \(m \mapsto xm\) 。因此我们有同构

\[
\begin{array}{l} \mathrm{H} _ {0} (x, \mathrm{M}) \to \mathrm{A} / x \mathrm{A} \otimes_ {\mathrm{A}} \mathrm{M} \leftarrow \mathrm{H} ^ {1} (x, \mathrm{M}), \\ \mathrm{H} _ {1} (x, \mathrm{M}) \to \mathrm{Hom} _ {\mathrm{A}} (\mathrm{A} / x \mathrm{A}, \mathrm{M}) \leftarrow \mathrm{H} ^ {0} (x, \mathrm{M}). \end{array}
\]

引理3. --- 设 K 是一个复形，使得 \(\mathbf{K}_i = 0\) 当 \(i \neq 0, 1\) ，并设 C 是一个复形， \(p\) 为一个整数。

\begin{enumerate}
\def\labelenumi{\alph{enumi})}
\tightlist
\item
  若 \(K\) 是平坦的，则对所有 \(p \in \mathbf{N}\) ，我们有正合列
\end{enumerate}

\[
0 \rightarrow \mathrm{H} _ {0} (\mathbf {K} \otimes_ {\mathbf {A}} \mathrm{H} _ {p} (\mathbf {C})) \rightarrow \mathrm{H} _ {p} (\mathbf {K} \otimes_ {\mathbf {A}} \mathbf {C}) \rightarrow \mathrm{H} _ {1} (\mathbf {K} \otimes_ {\mathbf {A}} \mathrm{H} _ {p - 1} (\mathbf {C})) \rightarrow 0.
\]

\begin{enumerate}
\def\labelenumi{\alph{enumi})}
\setcounter{enumi}{1}
\tightlist
\item
  若 K 是投射的，则对每个 \(p \in \mathbf{N}\) ，我们有正合列
\end{enumerate}

\[
\begin{array}{r l} 0 \to \mathrm{H} ^ {1} (\operatorname{Homgr} _ {\mathbf {A}} (\mathbf {K}, \mathrm{H} ^ {p - 1} (\mathbf {C}))) & \to \mathrm{H} ^ {p} (\operatorname{Homgr} _ {\mathbf {A}} (\mathbf {K}, \mathbf {C})) \\ & \to \mathrm{H} ^ {0} (\operatorname{Homgr} _ {\mathbf {A}} (\mathbf {K}, \mathrm{H} ^ {p} (\mathbf {C}))) \to 0. \end{array}
\]

我们来证明 \(a\) ， \(b\) ) 的证明类似。对每个 \(i\) ，记 \(\mathbf{K}_{(i)}\) 为复形 \(\mathbf{K}_i(-i)\) 。我们有一个复形的正合列，它作为 A-模的序列是分裂的

\[
0 \rightarrow \mathrm{K} _ {(0)} \xrightarrow {\alpha} \mathrm{K} \xrightarrow {\beta} \mathrm{K} _ {(1)} \rightarrow 0;
\]

序列

\[
0 \longrightarrow \mathrm{K} _ {(0)} \otimes_ {\mathrm{A}} \mathrm{C} \xrightarrow {\alpha \otimes 1} \mathrm{K} \otimes_ {\mathrm{A}} \mathrm{C} \xrightarrow {\beta \otimes 1} \mathrm{K} _ {(1)} \otimes_ {\mathrm{A}} \mathrm{C} \longrightarrow 0\tag{14}
\]

是正合的，并且由于 K 是平坦的，同态

\[
\gamma_ {0, p} (\mathbf {K} _ {(0)}, \mathbf {C}): \mathbf {K} _ {(0)} \otimes_ {\mathbf {A}} \mathrm{H} _ {p} (\mathbf {C}) \rightarrow \mathrm{H} _ {p} (\mathbf {K} _ {(0)} \otimes_ {\mathbf {A}} \mathbf {C})
\]

\[
\gamma_ {1, p - 1} (\mathbf {K} _ {(1)}, \mathbf {C}): \mathbf {K} _ {(1)} \otimes_ {\mathbf {A}} \mathrm{H} _ {p - 1} (\mathbf {C}) \rightarrow \mathrm{H} _ {p} (\mathbf {K} _ {(1)} \otimes_ {\mathbf {A}} \mathbf {C})
\]

是双射的（X，第 66 页，推论 2）。让我们计算连接同态 \(\partial (\alpha \otimes 1, \beta \otimes 1)\) ；根据定义，它把闭链 \(\sum a_i \otimes b_i\) 的类映到 \(\sum da_i \otimes b_i\) 的类，这意味着

\[
\partial (\alpha \otimes 1, \beta \otimes 1) \circ \gamma (\mathbf {K} _ {(1)}, \mathbf {C}) = \gamma (\mathbf {K} _ {(0)}, \mathbf {C}) \circ (d _ {\mathbf {K}} \otimes 1).
\]

因此，与 (14) 相关联的同调正合列具有如下形式：

\[
\begin{array}{r l} \mathrm{K} _ {1} \otimes \mathrm{H} _ {p} (\mathrm{C}) & \xrightarrow {d _ {\mathbf {k}} \otimes 1} \mathrm{K} _ {0} \otimes \mathrm{H} _ {p} (\mathrm{C}) \longrightarrow \mathrm{H} _ {p} (\mathrm{K} \otimes \mathrm{C}) \\ & \longrightarrow \mathrm{K} _ {1} \otimes \mathrm{H} _ {p - 1} (\mathrm{C}) \xrightarrow {d _ {\mathbf {k}} \otimes 1} \mathrm{K} _ {0} \otimes \mathrm{H} _ {p - 1} (\mathrm{C}), \end{array}
\]

由此得 a）。

将引理 3 的 a) 部分应用于复形 \(\mathbf{K}(u)\) ，并利用交换同构，我们得到：

命题 4. --- 对每个复形 C，我们有正合列。

\[
0 \rightarrow \mathrm{A} / x \mathrm{A} \otimes_ {\mathrm{A}} \mathrm{H} _ {p} (\mathrm{C}) \rightarrow \mathrm{H} _ {p} (x, \mathrm{C}) \rightarrow \operatorname{Hom} _ {\mathrm{A}} (\mathrm{A} / x \mathrm{A}, \mathrm{H} _ {p - 1} (\mathrm{C})) \rightarrow 0.
\]

推论 1. --- \(\mathbf{H}_{p}(x,\mathbf{C})=0\) 当且仅当比率为 \(x\) 的位似到 \(\mathbf{H}_{p}(\mathbf{C})\) 上是满射的，且比率为 \(x\) 的位似到 \(\mathbf{H}_{p-1}(\mathbf{C})\) 上是单射的。

推论 2. --- 设 \(x = (x_1, \ldots, x_n)\) 是 A 中元素的一个序列，M 是一个 A-模， \(x'\) 为序列 \((x_1, \ldots, x_{n-1})\) 。我们有正合列

\[
0 \rightarrow \mathrm{A} / x _ {n} \mathrm{A} \otimes_ {\mathrm{A}} \mathrm{H} _ {p} (x ^ {\prime}, \mathrm{M}) \rightarrow \mathrm{H} _ {p} (x, \mathrm{M}) \rightarrow \operatorname{Hom} _ {\mathrm{A}} (\mathrm{A} / x _ {n} \mathrm{A}, \mathrm{H} _ {p - 1} (x ^ {\prime}, \mathrm{M})) \rightarrow 0.
\]

推论 3. --- \(\mathrm{H}_{i}(\mathbf{x},\mathbf{M})\) 对所有 i \textgreater{} 0 均为零，当且仅当比率为 \(x_{n}\) 的位似到 \(\mathrm{H}_{i}(\mathbf{x}^{\prime},\mathbf{M})\) 上对 i \textgreater{} 0 是双射，并且比率为 \(x_{n}\) 的位似到 \(\mathrm{M}/(\mathbf{x}^{\prime})\) M 上是单射。

